%% F08 --- Trigonometry and the unit circle
%%
%% Every number in this program that the reader cannot do in their head comes
%% from code/f08_trigonometry.py and is pulled in with \val{}. The values a
%% reader must be able to read straight off the circle -- cos 0 = 1,
%% sin(pi/2) = 1, cos pi = -1, cos^2 + sin^2 = 1 -- are written inline,
%% because looking them up would defeat the program.
%%
%% THIS PROGRAM IS NOT A TRIGONOMETRY COURSE. There is no sine rule, no cosine
%% rule, no solving of triangles and no list of identities to memorise. The
%% book's scope statement excludes all of it. What is here is the one picture
%% the rest of the book needs -- a point going round a circle -- and the two
%% things that picture buys.
%%
%% THE TWO PAYOFFS, and both are IDENTITIES asserted over a sweep in the script
%% rather than shown at a point:
%%   cosine similarity IS the cosine of an angle, checked against the angle
%%     computed independently, worst error 7.8e-16 over 3481 pairs;
%%   R(a)q . R(b)k depends only on b - a, worst error 3.6e-15 over 1600 pairs,
%%     which is the whole of why rotating a query and a key by an angle
%%     proportional to position encodes RELATIVE position.
%% The second is the sentence the brief asks for: it makes "rotate the query
%% and key" a claim the reader can check rather than an incantation.
%%
%% THE TRAP IS THE UNITS and it is computed rather than described.
%% math.sin(90) is 0.8940, because 90 is read as radians and is 14.32 full
%% turns. Nothing raises and nothing warns; the number just looks like nothing.
%%
%% Twin: programs/pl/F08-trigonometry.tex. Frame for frame, macro for macro,
%% number for number; tools/parity.py compares them.

\program{Trigonometry and the unit circle}
\label{prog:F08}
\index{trigonometry!unit circle}
\index{cosine similarity!is a cosine}

\begin{outcomes}
\outcome{1--6}{read $\cos$ and $\sin$ off a point going round a circle, work in
  radians, and get the one identity this book needs straight from the picture}
\outcome{7--10}{sketch both curves from the circle, and say what their period,
  their range and their symmetry are without recalling any of them}
\outcome{11--15}{apply Program~\ref{prog:F05}'s four moves to a wave, and say
  what each of the four numbers in $A\sin(B\theta + C) + D$ controls}
\outcome{16--23}{compute a cosine similarity, say what it measures and what it
  ignores, and say what it does \emph{not} tell you}
\outcome{24--30}{rotate a point in the plane, say what a rotation leaves
  unchanged, and say what that buys when two vectors are turned by different
  amounts}
\end{outcomes}

\begin{quiz}
\item A point sits on a circle of radius $1$, a quarter turn anticlockwise
  from the right-hand side. What are its coordinates? \teachesat{2--3}
  \answerto{$(0, 1)$, so $\cos$ of a quarter turn is $0$ and $\sin$ of it
  is $1$.}
\item How many radians are there in a full turn? \teachesatone{4}
  \answerto{$2\pi$, because a radian is defined so that the angle equals the
  distance walked round a circle of radius $1$.}
\item What does \code{math.sin(90)} return? \teachesat{4--5}
  \answerto{$\val{f08.sin90rad}$. The argument is read as radians, and $90$
  radians is $\val{f08.turns90}$ full turns.}
\item A point $(\cos\theta, \sin\theta)$ sits on a circle of radius $1$. What
  equation does that give you? \teachesat{5--6}
  \answerto{$\cos^{2}\theta + \sin^{2}\theta = 1$, because that equation is the
  circle of radius $1$ with the names of the coordinates substituted in.}
\item What is the period of $\sin$, and what is its range? \teachesat{7--8}
  \answerto{$2\pi$ and $\intcc{-1}{1}$. Both are properties of going round a
  circle of radius $1$ and returning to where you began.}
\item Write $\cos$ in terms of $\sin$. \teachesat{8--9}
  \answerto{$\cos\theta = \sin\!\left(\theta + \frac{\pi}{2}\right)$: the same
  curve shifted by a quarter turn.}
\item In $A\sin(B\theta + C) + D$, which letter sets how many waves fit into a
  given stretch? \teachesat{12--13}
  \answerto{$B$. It is inside the bracket, so it squashes the wave
  horizontally, exactly as Program~\ref{prog:F05} said.}
\item Two vectors point in exactly the same direction but one is ten times
  longer. What is their cosine similarity? \teachesat{16--17}
  \answerto{$1$. The lengths divide out, which is the whole point of the
  measure.}
\item What is the cosine similarity of two vectors at right angles?
  \teachesat{17--18}
  \answerto{$0$, because $\cos$ of a quarter turn is $0$.}
\item A cosine similarity of $\num{0.99}$ is reported. What angle is that?
  \teachesat{18--19}
  \answerto{About $\val{f08.ang.hi}$ degrees. Cosine is flat near zero, so a
  similarity close to $1$ still leaves a visible angle.}
\item Does a high cosine similarity between two embeddings mean the two things
  are similar in meaning? \teachesat{20--21}
  \answerto{Not on its own. It means the vectors point the same way, and
  whether that corresponds to meaning is a property of how they were trained.}
\item Write down what a rotation by $\theta$ does to the point $(x, y)$.
  \teachesatone{24}
  \answerto{It sends it to
  $(x\cos\theta - y\sin\theta,\; x\sin\theta + y\cos\theta)$.}
\item What does a rotation do to the length of a vector? \teachesat{25--26}
  \answerto{Nothing. A rotation moves a point round a circle centred on the
  origin, and every point of that circle is the same distance out.}
\item Two vectors are both rotated by the same angle. What happens to their
  dot product? \teachesat{26--27}
  \answerto{Nothing. The angle between them has not changed and neither has
  either length, and the dot product depends only on those three things.}
\item A query at position $m$ and a key at position $n$ are each rotated by an
  angle proportional to their position. What does their dot product depend on?
  \teachesat{27--28}
  \answerto{Only $n - m$. Both absolute positions cancel, which is why the
  scheme encodes relative position.}
\end{quiz}

% ============================================================
\section{A point going round a circle}
\label{sec:F08-circle}
\index{radian!as arc length}

\begin{fr}
Everything in this program comes from one picture, and the picture is a point
walking anticlockwise round a circle of radius $1$ centred on the origin,
starting from the right-hand side.

At any moment the point is somewhere, and \enquote{somewhere} in the plane
means two numbers. Those two numbers have names:

\begin{center}
\begin{tabular}{ll}
\toprule
$\cos\theta$ & how far along the point is, across \\
$\sin\theta$ & how far up it is \\
\bottomrule
\end{tabular}
\end{center}

That is the whole definition. Not a ratio of sides, not something about
right-angled triangles \dash{} the coordinates of a point on a circle, with
$\theta$ saying how far round it has gone.

The point starts at the right-hand side. What are its coordinates there, and
therefore what are $\cos 0$ and $\sin 0$?
\nextframe
\end{fr}

\begin{fr}
\begin{ansblock}
$(1, 0)$, so $\cos 0 = 1$ and $\sin 0 = 0$.
\end{ansblock}

Now walk it a quarter of the way round, anticlockwise, so that it sits at the
top of the circle. What are its coordinates?
\nextframe
\end{fr}

\begin{fr}
\ans{$(0, 1)$}
Directly above the origin, one unit up and none across. So $\cos$ of a quarter
turn is $0$ and $\sin$ of a quarter turn is $1$.

Carry on and the rest come for free, with no table to memorise:

\begin{center}
\begin{tabular}{lcc}
\toprule
Position & $\cos$ & $\sin$ \\
\midrule
right & $1$ & $0$ \\
top & $0$ & $1$ \\
left & $-1$ & $0$ \\
bottom & $0$ & $-1$ \\
\bottomrule
\end{tabular}
\end{center}

Read the columns rather than the rows and you have the two curves already: one
starts at $1$ and falls, the other starts at $0$ and rises.

One thing is still missing, and it is the thing everybody gets wrong in code.
How should $\theta$ be measured?
\nextframe
\end{fr}

\begin{fr}
\begin{ansblock}
In radians: $\theta$ is the \emph{distance the point has walked} round the rim.
\end{ansblock}

That is the whole definition of a radian, and it is the reason the unit exists.
The circle has radius $1$, so its circumference is $2\pi$, so \textbf{a full
turn is $2\pi$ radians} and a quarter turn is $\frac{\pi}{2}$. The angle and
the arc length are the same number, which is what makes every later formula
come out without a conversion factor sitting in it.

Degrees are a human convention with no mathematics behind them \dash{} $360$
was chosen because it divides conveniently \dash{} and every library in your
working life takes radians.

Which brings us to a question to answer quickly. \textbf{What does
\code{math.sin(90)} return?}
\dotline
\nextframe
\end{fr}

\begin{fr}
\ans{$\val{f08.sin90rad}$}
Not $1$.

\begin{trapbox}
\textbf{The argument is read as radians, and $90$ radians is
$\val{f08.turns90}$ full turns.} The point has gone round more than fourteen
times and stopped somewhere unremarkable, and $\val{f08.sin90rad}$ is where it
happened to land.

What makes this expensive is everything that does \emph{not} happen. Nothing
raises, nothing warns, and the value is a perfectly ordinary number between
$-1$ and $1$ \dash{} so it flows into whatever comes next and the first sign
of trouble is a result that is merely wrong.

The habit that prevents it: \textbf{write \code{math.radians(90)} whenever the
number in your head is in degrees}, even when you are sure, because the cost
of the conversion is nothing and the cost of the mistake is a debugging
session that starts nowhere near the bug.
\end{trapbox}

One identity is worth having, and it needs no proof. The point $(\cos\theta,
\sin\theta)$ is on a circle of radius $1$. What does that tell you?
\nextframe
\end{fr}

\begin{fr}
\begin{ansblock}
\[ \cos^{2}\theta + \sin^{2}\theta = 1 \]
\end{ansblock}

Because that \emph{is} the equation of a circle of radius $1$: a point $(x, y)$
is on it exactly when $x^{2} + y^{2} = 1$, which is Pythagoras applied to the
horizontal and vertical distances. Substitute the names and you are done.

So the most-quoted identity in trigonometry is not a fact to memorise, it is
the definition read back. \code{code/f08\_trigonometry.py} checks it at
fourteen hundred angles, with a worst disagreement of $\val{f08.pyth.err}$,
which is as close to exactly $1$ as binary64 gets.

\begin{note}
This is the only trigonometric identity this book uses, and it is here because
it falls out of the picture. There is no sine rule here, no cosine rule, no
solving of triangles and no list of formulas to learn \dash{} the book's scope
statement excludes all of it, and nothing in the rest of the book needs it.
What the rest of the book needs is the circle, and you now have it.
\end{note}
\end{fr}

\mermaidfig{f08-unit-circle}{The definition, in one line: the angle is how far
the point has walked, and the two coordinates are the two
functions.}{sine and cosine from the circle}

% ============================================================
\section{The two curves}
\label{sec:F08-curves}
\index{trigonometry!period and range}

\begin{fr}
Unroll the walk. Put $\theta$ along the horizontal axis and plot the height of
the point, and you get the sine curve; plot how far across it is instead, and
you get the cosine curve.

Everything about both curves is a fact about walking round a circle:

\begin{itemize}
\item they never leave $\intcc{-1}{1}$, because the point never leaves a
  circle of radius $1$;
\item they repeat every $2\pi$, because that is a full turn and the point is
  back where it started;
\item neither has a corner anywhere, because the walk has no corners.
\end{itemize}

None of that needs remembering. What is the period of $\sin$, and what is its
range?
\nextframe
\end{fr}

\begin{fr}
\begin{ansblock}
$2\pi$ and $\intcc{-1}{1}$.
\end{ansblock}

And the same for $\cos$, for the same reason.

Now the symmetry, which Program~\ref{prog:F05} gave you the words for. Walking
backwards round the circle by $\theta$ puts the point at the mirror image, in
the horizontal axis, of where walking forwards puts it. So the across-value is
unchanged and the up-value flips sign:
\[ \cos(-\theta) = \cos\theta \qquad \sin(-\theta) = -\sin\theta \]

$\cos$ is \emph{even} and $\sin$ is \emph{odd}, in exactly
Program~\ref{prog:F05}'s sense, and again the picture is the proof.

The two curves look like the same shape moved. Write $\cos$ in terms of $\sin$.
\nextframe
\end{fr}

\begin{fr}
\begin{ansblock}
\[ \cos\theta = \sin\!\left(\theta + \frac{\pi}{2}\right) \]
\end{ansblock}

Cosine is sine started a quarter turn early. Checked at fourteen hundred
angles with a worst disagreement of $\val{f08.shift.err}$.

That is Program~\ref{prog:F05}'s horizontal shift, and it has the direction
that program warned about: the $+\frac{\pi}{2}$ is inside the bracket, so the
graph moves the \emph{other} way \dash{} sine shifted a quarter turn to the
left is cosine.

\textbf{So there is one curve here, not two}, in the same way that
Program~\ref{prog:F07} found one curve behind the logistic and $\tanh$. The
book uses both names because both are conventional, not because there are two
things to learn.

Which of $\cos$ and $\sin$ is $1$ at $\theta = 0$, and which is $0$?
\nextframe
\end{fr}

\begin{fr}
\begin{ansblock}
$\cos 0 = 1$ and $\sin 0 = 0$.
\end{ansblock}

Which is the only thing you need to tell them apart, and it comes straight from
where the point starts.

\begin{aibox}
Sine and cosine both appear in a positional encoding, at alternating
coordinates, and the reason is this frame: at position zero the sine
coordinates read $0$ and the cosine coordinates read $1$, so the pattern at
position zero is not all zeros and the encoding can distinguish it from a
position that is genuinely absent.

Using one function for every coordinate would give a first position encoded as
a vector of zeros, which is the same thing a padded slot looks like.
\end{aibox}
\end{fr}

% ============================================================
\section{Four moves on a wave}
\label{sec:F08-moves}
\index{wave!amplitude, frequency, phase}

\begin{fr}
A wave in the wild is never the bare $\sin\theta$. It is
\[ y = A\sin(B\theta + C) + D \]
and all four letters are Program~\ref{prog:F05}'s four moves, applied to a
curve you now know.

Two are outside the bracket and act on the output, so they read as written:
$A$ stretches the curve vertically and $D$ slides it up. Two are inside and
act on the input, so they read backwards: $B$ squashes it horizontally and $C$
slides it sideways, in the direction opposite to its sign.

Give the range of $y = 3\sin\theta + 2$.
\nextframe
\end{fr}

\begin{fr}
\ans{From $-1$ to $5$}
Because $\sin$ runs from $-1$ to $1$, tripling gives $-3$ to $3$, and adding
$2$ gives $-1$ to $5$. $A$ is the \textbf{amplitude} and $D$ is the offset,
and they are the two easy ones.

Now the inside pair, which are the ones with names worth knowing. How many
complete waves does $\sin(2\theta)$ fit into one full turn of $\theta$?
\nextframe
\end{fr}

\begin{fr}
\ans{Two}
Because $2\theta$ reaches $2\pi$ when $\theta$ reaches $\pi$, so the wave has
completed a cycle in half the space. $B$ is the \textbf{frequency}, and the
wavelength is $\frac{2\pi}{B}$ \dash{} bigger $B$, shorter wave, which is
Program~\ref{prog:F05}'s inside-the-bracket squash with a name attached.

$C$ is the \textbf{phase}: it slides the wave sideways by $\frac{C}{B}$, and
the division by $B$ is there for the reason Program~\ref{prog:F05} gave when
it put the logistic's crossing point at $-\frac{b}{w}$ rather than at $b$. The
shift is measured after the squash.

What is the wavelength of $\sin(\num{0.5}\,\theta)$?
\nextframe
\end{fr}

\begin{fr}
\ans{$4\pi$}
From $\frac{2\pi}{\num{0.5}}$. A small frequency is a long wave.

\begin{aibox}
A sinusoidal positional encoding is this expression with many different values
of $B$ at once: one pair of coordinates per frequency, and the frequencies
spread over an enormous range. In the original transformer's encoding at
$d_{\text{model}} = 512$, the wavelengths run from about
$\val{f08.pe.first}$ positions at the fastest coordinate to about
$\val{f08.pe.last}$ at the slowest, with $\val{f08.pe.mid}$ in the middle.

That spread is the design. A wave of wavelength $\val{f08.pe.first}$ tells
neighbouring positions apart and says nothing about positions a thousand
tokens away, because it has gone round hundreds of times in between. A wave of
wavelength $\val{f08.pe.last}$ has barely started moving over the whole
sequence and so cannot distinguish neighbours, but it does distinguish the
beginning from the end.

Neither wave can do the other's job. Together they carry position at every
scale at once, which is what a single frequency could not.
\end{aibox}
\end{fr}

\begin{fr}
One caution about frequency, because it is the mirror of the units trap and
just as easy to acquire.

\textbf{A wavelength is a number of positions, not a number of radians.} The
wave $\sin(B\theta)$ completes a cycle when $B\theta$ advances by $2\pi$, so
$\theta$ advances by $\frac{2\pi}{B}$ \dash{} and it is that second number,
not $2\pi$, that tells you how many tokens apart two positions must be before
the coordinate has come back round to look the same.

Two positions $\frac{2\pi}{B}$ apart are indistinguishable in that one
coordinate. They are told apart by the slower coordinates, and that is the
whole reason there is more than one of them.
\end{fr}

% ============================================================
\section{Cosine similarity is a cosine}
\label{sec:F08-cosine-similarity}
\index{cosine similarity!what it ignores}

\begin{fr}
Here is the formula you have used and possibly never unpacked. For two
vectors $a$ and $b$,
\[ \cos\theta = \frac{a \cdot b}{\lVert a \rVert\, \lVert b \rVert} \]
where $a \cdot b$ is Program~\ref{prog:F04}'s $\sum_i a_i b_i$ and
$\lVert a \rVert$ is the length.

\textbf{The $\theta$ on the left is an angle, and this is an equality, not an
analogy.} Cosine similarity is not \emph{like} a cosine and it is not
\emph{named after} one. It is the cosine of the angle between the two vectors,
and \code{code/f08\_trigonometry.py} checks that against the angle computed
independently over three and a half thousand pairs, worst disagreement
$\val{f08.cossim.err}$.

Two vectors point in exactly the same direction and one is ten times longer
than the other. What is their cosine similarity?
\nextframe
\end{fr}

\begin{fr}
\ans{$1$}
The lengths appear in the denominator precisely so that they cancel. That is
the design of the measure and it is what distinguishes it from the bare dot
product: \textbf{cosine similarity is about direction and knows nothing about
magnitude.}

Whether that is what you want is a separate question, and it is worth asking
rather than assuming. If the length of an embedding carries information
\dash{} frequency, confidence, how much text went into it \dash{} then cosine
similarity is throwing it away.

Work out the cosine similarity of $(3, 4)$ and $(-4, 3)$.
\nextframe
\end{fr}

\begin{fr}
\ans{$\val{f08.cos.right}$}
Because the dot product is $-12 + 12 = 0$. And a cosine similarity of $0$ means
an angle of $\val{f08.ang.right}$ degrees, since $\cos$ of a quarter turn is
zero \dash{} the two vectors are at right angles.

The three values worth being able to read at sight:

\begin{center}
\begin{tabular}{lcc}
\toprule
& similarity & angle \\
\midrule
same direction & $\val{f08.cos.same}$ & $\val{f08.ang.same}$° \\
at right angles & $\val{f08.cos.right}$ & $\val{f08.ang.right}$° \\
opposite & $\val{f08.cos.opposite}$ & $\val{f08.ang.opposite}$° \\
\bottomrule
\end{tabular}
\end{center}

Now a number you will actually meet. A retrieval system reports a cosine
similarity of $\val{f08.cos.hi}$ between two embeddings. What angle is that?
\nextframe
\end{fr}

\begin{fr}
\ans{About $\val{f08.ang.hi}$ degrees}
Which is a visible angle, not a rounding error.

\begin{trapbox}
\textbf{Cosine is flat near zero, so similarities close to $1$ compress a wide
range of angles into a narrow range of numbers.} That is
Program~\ref{prog:F07}'s saturation in a different function: $\cos$ has its
turning point at $\theta = 0$, so the first few degrees of separation cost
almost nothing in similarity.

The consequence is one people trip over when tuning a threshold. Moving a
cutoff from $\val{f08.cos.hi}$ to $\val{f08.cos.lo}$ sounds like a small
adjustment because the numbers are close, and it widens the accepted angle
from about $\val{f08.ang.hi}$ degrees to about $\val{f08.ang.lo}$. The
similarity scale is not linear in the thing you care about, and it is least
linear exactly where you are most likely to set the threshold.
\end{trapbox}
\end{fr}

\begin{fr}
And now the harder question, which is not about mathematics at all.

Two document embeddings have a cosine similarity of $\num{0.95}$. Does that
mean the two documents are similar in meaning?
\dotline
\nextframe
\end{fr}

\begin{fr}
\ans{Not on its own}
It means the two vectors point in nearly the same direction. Whether that
corresponds to meaning is a property of \emph{how the vectors were made}, and
it is a claim about the training procedure rather than about the geometry.

\begin{trapbox}
\textbf{The mathematics is exact and the interpretation is empirical, and it
is easy to let the exactness of the first launder the second.} The cosine is
the cosine of an angle, to fifteen decimal places. That the angle tracks
meaning is a fact about a model, holds better for some pairs than others, and
has to be measured on the task rather than assumed from the formula.

Two failure modes are reported often enough to expect: embeddings from many
models are squeezed into a narrow cone, so \emph{everything} is similar to
everything and a raw score of $\num{0.8}$ may be unremarkable; and a
similarity can be driven by shared topic, shared style, shared length or
shared language rather than by what the text says. Neither is measured in this
book. Program~\ref{prog:P05} does measure the baseline they are deviations
from \dash{} in high dimension two unrelated vectors are almost always nearly
orthogonal \dash{} and Program~\ref{prog:P34} says how to find out whether a
score means anything on your data.

What you may take from this program is the geometry, which is exact. What you
may not take is that the geometry is the meaning.
\end{trapbox}
\end{fr}

\mermaidfig{f08-cosine-angle}{Dividing by both lengths removes them, and what
survives is an angle.}{cosine similarity unpacked}

\begin{fr}
One last property, because it is where the next section starts. Cosine
similarity depends on the two vectors only through the angle between them.

That means anything you do to both vectors that leaves the angle alone leaves
the similarity alone. Scaling both is the obvious example and you have already
seen it.

Name another operation that would leave the angle between two vectors
unchanged.
\nextframe
\end{fr}

\begin{fr}
\begin{ansblock}
Turning both of them by the same amount.
\end{ansblock}

Rotate the whole picture and the vectors move but the angle between them does
not, because the angle is a relationship between the two rather than a
property of either. That is obvious from the picture and it is the entire
subject of section 5, where it turns out to be worth a great deal more than it
looks.
\end{fr}

% ============================================================
\section{Rotation, and what it is for}
\label{sec:F08-rotation}
\index{rotation!preserves the dot product}

\begin{fr}
To turn a point anticlockwise by an angle $\theta$, send $(x, y)$ to
\[ (x\cos\theta - y\sin\theta,\quad x\sin\theta + y\cos\theta) \]

You do not need to derive that here, but you should check it once on a case you
can see. Take the point $(1, 0)$, on the right-hand side of the circle, and
rotate it by a quarter turn.

Where does it go?
\nextframe
\end{fr}

\begin{fr}
\ans{To $(0, 1)$}
Because $\cos$ of a quarter turn is $0$ and $\sin$ of it is $1$, so the formula
gives $(1 \times 0 - 0 \times 1,\; 1 \times 1 + 0 \times 0) = (0, 1)$. The
point on the right has moved to the top, which is what a quarter turn
anticlockwise does.

What does a rotation do to the \emph{length} of a vector?
\nextframe
\end{fr}

\begin{fr}
\ans{Nothing}
A rotation moves a point round a circle centred on the origin, and every point
of that circle is the same distance from the origin. Checked over six hundred
angles, worst change in length $\val{f08.rot.len.err}$.

That is worth stating as the property rather than the picture, because it is
the property everything else rests on: \textbf{a rotation changes where things
are and not how big they are.}

Now take two vectors and rotate \emph{both} by the same angle. What happens to
their dot product?
\nextframe
\end{fr}

\begin{fr}
\ans{Nothing}
Worst change $\val{f08.rot.same.err}$ over six hundred angles.

And you can see why without any algebra, from the previous section. The dot
product is
\[ a \cdot b = \lVert a \rVert\, \lVert b \rVert \cos\theta \]
by rearranging the cosine-similarity formula. A rotation leaves both lengths
alone, and it leaves the angle between the two vectors alone because it turns
them both by the same amount. All three ingredients survive, so the product
does.

Now the step that matters. Rotate $a$ by an angle $\alpha$ and $b$ by a
\emph{different} angle $\beta$. What can their dot product depend on?
\nextframe
\end{fr}

\begin{fr}
\begin{ansblock}
Only on $\beta - \alpha$.
\end{ansblock}

Turn the whole picture back by $\alpha$. That is a rotation of both vectors by
the same amount, so by the previous frame it changes nothing \dash{} and it
leaves $a$ where it started and $b$ rotated by $\beta - \alpha$. So
\[ R(\alpha)a \cdot R(\beta)b = a \cdot R(\beta - \alpha)b \]
and the two absolute angles have gone. Checked over sixteen hundred pairs of
angles, worst disagreement $\val{f08.rot.rel.err}$.

\textbf{That is the whole of it}, and it is the sentence this program exists to
make sayable.
\end{fr}

\begin{fr}
Here it is with the vocabulary attached.

Give each position in a sequence an angle proportional to it. Rotate the query
at position $m$ by its angle and the key at position $n$ by its angle. Then
their dot product \dash{} which is the attention score \dash{} depends on the
two vectors and on $n - m$, and on nothing else.

Neither token knows where it is in absolute terms. The score knows only how far
apart they are.

The script does it on one pair. At positions $1$ and $2$ a query and a key
one position apart score $\val{f08.rot.score.near}$; at positions $9$ and
$10$ the same two vectors score $\val{f08.rot.score.far}$.

\begin{aibox}
That is rotary position embedding, and \enquote{rotate the query and the key}
is now a sentence rather than an incantation: it is the only way to put
position into a dot product such that the answer depends on the difference of
the positions and not on the positions themselves.

Two consequences fall straight out of the identity and are worth having. A
sequence shifted bodily along produces identical attention scores, because
every difference is unchanged \dash{} which is the property you want and is
not shared by adding a positional vector to the embedding. And the scheme
needs no learned parameters at all, because the rotation is determined by the
position and the frequency schedule and there is nothing to fit.

What it does not do is extend for free beyond the range of angles the model
was trained on. The identity holds at every angle; whether the model has seen
that part of the circle is a different question, and it is the one the
extrapolation literature is about.
\end{aibox}
\end{fr}

\mermaidfig{f08-rotate-both}{Two rotations by two positions, and a dot product
that remembers only the gap between them.}{why rotation encodes a difference}

\begin{fr}
One closing observation, which ties the section to the one before it.

Section 4 said cosine similarity is unchanged by rotating both vectors, and
called that obvious. Section 5 said an attention score is unchanged by rotating
both vectors, and made it the mechanism behind a real technique.

They are the same statement. What changed is that the rotation is no longer
something you do to the whole picture at once \dash{} it is something you do to
each vector by a different amount, chosen so that what survives is exactly the
quantity you wanted to encode.

\textbf{Rotation is the operation that changes where things are without
changing how they relate}, and that is precisely why it is the right tool for
putting position into a model that must not otherwise notice it.
\end{fr}

% ============================================================
\begin{summarybox}
\sumitem{1--3}{\textbf{$\cos\theta$ and $\sin\theta$ are the two coordinates
  of a point on a circle of radius $1$}, after it has walked $\theta$
  anticlockwise from the right-hand side. Not a ratio of sides and nothing to
  do with triangles \dash{} the whole subject is one picture, and the four
  quarter-turn values are read off it rather than memorised.}
\sumitem{4--5}{\textbf{An angle is a distance walked, which is what a radian
  is}, so a full turn is $2\pi$ and no conversion factor survives into any
  later formula. \code{math.sin(90)} returns $\val{f08.sin90rad}$, because
  $90$ radians is $\val{f08.turns90}$ full turns \dash{} and nothing raises,
  nothing warns, and the answer is an ordinary number.}
\sumitem{5--6}{$\cos^{2}\theta + \sin^{2}\theta = 1$ is
  \textbf{the equation of the circle with the names substituted in}, so it
  needs no proof and no memorising. Checked to $\val{f08.pyth.err}$. It is
  the only trigonometric identity this book uses.}
\sumitem{7--8}{Both curves have period $2\pi$ and range $\intcc{-1}{1}$
  because the point returns to where it started and never leaves the circle.
  $\cos$ is even and $\sin$ is odd, which is walking backwards mirrored in the
  horizontal axis.}
\sumitem{9--10}{\textbf{$\cos\theta = \sin(\theta + \frac{\pi}{2})$}, checked
  to $\val{f08.shift.err}$: one curve, two names, and the $+$ inside the
  bracket moves the graph the other way, as Program~\ref{prog:F05} warned. At
  $\theta = 0$ they read $1$ and $0$, which is the only thing that tells them
  apart.}
\sumitem{11--13}{In $A\sin(B\theta + C) + D$ the four letters are
  Program~\ref{prog:F05}'s four moves: $A$ the amplitude and $D$ the offset
  act on the output and read as written; $B$ the frequency and $C$ the phase
  act on the input and read backwards. The wavelength is $\frac{2\pi}{B}$ and
  the sideways slide is $\frac{C}{B}$, divided by $B$ because the shift is
  measured after the squash.}
\sumitem{14--15}{\textbf{A wavelength is a number of positions, not a number
  of radians.} A sinusoidal positional encoding is one wave per frequency over
  an enormous spread \dash{} about $\val{f08.pe.first}$ positions at the
  fastest and $\val{f08.pe.last}$ at the slowest \dash{} because neither end
  can do the other's job.}
\sumitem{16--17}{\textbf{Cosine similarity is the cosine of the angle between
  two vectors}, not an analogy and not a name. Checked against the angle
  computed independently, worst disagreement $\val{f08.cossim.err}$. The two
  lengths are in the denominator so that they cancel, which means the measure
  is about direction and knows nothing about magnitude \dash{} and whether
  that is what you want is a question worth asking.}
\sumitem{18--19}{$\val{f08.cos.same}$, $\val{f08.cos.right}$ and
  $\val{f08.cos.opposite}$ are the same direction, right angles and opposite.
  \textbf{Cosine is flat near zero}, so a similarity of $\val{f08.cos.hi}$ is
  still about $\val{f08.ang.hi}$ degrees, and dropping the cutoff to
  $\val{f08.cos.lo}$ admits $\val{f08.ang.lo}$ \dash{} the scale is least
  linear exactly where thresholds get set.}
\sumitem{20--21}{\textbf{The mathematics is exact and the interpretation is
  empirical.} A high similarity says the vectors point the same way; that this
  tracks meaning is a claim about how they were trained, and it has to be
  measured on the task.}
\sumitem{22--23}{Cosine similarity depends on the two vectors only through the
  angle between them, so scaling both leaves it alone \dash{} and so does
  turning both by the same amount.}
\sumitem{24--26}{A rotation sends $(x, y)$ to
  $(x\cos\theta - y\sin\theta,\ x\sin\theta + y\cos\theta)$ and
  \textbf{changes where things are without changing how big they are}: worst
  change in length $\val{f08.rot.len.err}$ over six hundred angles.}
\sumitem{27--28}{Rotating \emph{both} vectors by the same angle leaves the dot
  product alone, to $\val{f08.rot.same.err}$, because all three of its
  ingredients survive. So rotating them by \emph{different} angles $\alpha$
  and $\beta$ leaves a dot product that can depend only on $\beta - \alpha$
  \dash{} checked to $\val{f08.rot.rel.err}$.}
\sumitem{29--30}{\textbf{That identity is rotary position embedding.} Turn the
  query and the key each by an angle proportional to its position and the
  attention score depends on the gap and not on the two positions: the same
  pair scores $\val{f08.rot.score.near}$ at positions $1$ and $2$ and
  $\val{f08.rot.score.far}$ at $9$ and $10$. It needs no learned parameters,
  and it does not by itself extend past the angles the model was trained on.}
\end{summarybox}

\canyou

% ============================================================
\begin{testexercises}
\item A point sits at the bottom of the unit circle. Give $\cos$ and $\sin$
  there.
  \answerto{$0$ and $-1$. Directly below the origin: none across, one unit
  down.}
\item How many radians is a half turn, and how many is a sixth of a turn?
  \answerto{$\pi$ and $\frac{\pi}{3}$, from $2\pi$ divided by $2$ and by $6$.}
\item A colleague writes \code{math.cos(180)} meaning the far side of the
  circle. What is wrong, what should it be, and what will go wrong at run
  time?
  \answerto{The argument is radians, so it should be
  \begin{center}\code{math.cos(math.radians(180))}\end{center}
  which is $-1$. At run time nothing goes wrong \dash{} it returns an ordinary
  number in $\intcc{-1}{1}$ and the result is merely wrong.}
\item Without looking anything up, give the range of $\cos$ and say why.
  \answerto{$\intcc{-1}{1}$, because $\cos$ is a coordinate of a point on a
  circle of radius $1$ and no coordinate of such a point can exceed $1$ in
  size.}
\item Give the amplitude, wavelength and offset of
  $y = 4\sin(\num{0.25}\,\theta) - 1$.
  \answerto{Amplitude $4$, wavelength $\frac{2\pi}{\num{0.25}} = 8\pi$,
  offset $-1$. Its range is $-5$ to $3$.}
\item Two coordinates of a positional encoding have wavelengths $10$ and
  $\num{10000}$. Which one distinguishes positions $3$ and $5$, and which one
  distinguishes positions $3$ and $\num{3000}$?
  \answerto{The wavelength-$10$ coordinate separates $3$ from $5$ and has gone
  round hundreds of times by position $\num{3000}$; the wavelength-$\num{10000}$
  one has barely moved between $3$ and $5$ and separates $3$ from $\num{3000}$.
  That is why there is more than one.}
\item Compute the cosine similarity of $(1, 0)$ and $(1, 1)$, and give the
  angle.
  \answerto{$\frac{1}{\sqrt{2}} \approx \num{0.7071}$, which is $45$ degrees.}
\item Two embeddings have cosine similarity $\num{1.0}$ but very different
  lengths. What do you know about them?
  \answerto{That one is a positive multiple of the other. Cosine similarity
  cannot see the lengths at all, so it cannot distinguish a vector from twice
  itself.}
\item A threshold is moved from $\num{0.98}$ to $\num{0.99}$ and far fewer
  results pass. Explain in terms of angles.
  \answerto{The accepted angle narrows from about $\val{f08.ang.lo}$ degrees
  to about $\val{f08.ang.hi}$. Cosine is flat near zero, so a small change in
  the number is a large change in the angle it admits.}
\item Rotate $(0, 1)$ by a quarter turn anticlockwise using the formula.
  \answerto{$(0 \times 0 - 1 \times 1,\ 0 \times 1 + 1 \times 0) = (-1, 0)$.
  The top of the circle moves to the left, which is what a quarter turn does.}
\item Two vectors of lengths $3$ and $5$ have a dot product of $\num{7.5}$.
  Both are rotated by $\num{1.2}$ radians. What is their dot product now?
  \answerto{$\num{7.5}$. A rotation of both changes neither length nor the
  angle between them, and the dot product depends on nothing else.}
\item A query at position $4$ and a key at position $7$ are rotated by angles
  proportional to their positions. Which other pair of positions gives the
  same score for the same two vectors?
  \answerto{Any pair three apart \dash{} $1$ and $4$, $100$ and $103$. The
  score depends only on the difference.}
\item Why can a rotation not be used to make one vector longer than another?
  \answerto{Because a rotation preserves length. It moves points round circles
  centred on the origin, and every point of one circle is the same distance
  out.}
\item Name the one trigonometric identity this program uses, and say where it
  comes from.
  \answerto{$\cos^{2}\theta + \sin^{2}\theta = 1$, which is the equation of
  the unit circle with $\cos$ and $\sin$ substituted for the coordinates.}
\end{testexercises}

% ============================================================
\begin{furtherproblems}
\item Show that $\sin(\theta + 2\pi) = \sin\theta$ using the circle rather
  than any identity.
  \answerto{Adding $2\pi$ walks the point one full turn round the rim, which
  returns it to exactly the same place. Both coordinates are therefore
  unchanged, and $\sin$ is one of them.}
\item A wave is described as having \enquote{frequency $3$}. Give its
  wavelength, and say why the two numbers move in opposite directions.
  \answerto{$\frac{2\pi}{3}$. $B$ multiplies the input, so a bigger $B$ makes
  the input reach $2\pi$ sooner \dash{} which is a shorter wave. It is
  Program~\ref{prog:F05}'s inside-the-bracket squash.}
\item Two vectors have cosine similarity $-\num{0.5}$. What angle is that, and
  what does it say about their directions?
  \answerto{$120$ degrees. They point more away from each other than towards
  each other, but not oppositely.}
\item A model's embeddings all lie within $20$ degrees of one another. What is
  the smallest cosine similarity any two of them can have, and why is that a
  problem for a threshold?
  \answerto{$\cos 20^{\circ} \approx \num{0.94}$. Every pair scores at least that, so
  a threshold of $\num{0.9}$ admits everything and the number carries no
  information about the pair.}
\item Prove that rotating both vectors by the same angle leaves their cosine
  similarity unchanged, given only that a rotation preserves length and the
  dot product.
  \answerto{The similarity is the dot product divided by the two lengths.
  A rotation of both leaves the numerator alone by assumption and each
  denominator factor alone by assumption, so it leaves the quotient alone.}
\item In $R(\alpha)a \cdot R(\beta)b = a \cdot R(\beta - \alpha)b$, what is the
  role of the assumption that the two rotations are about the same origin?
  \answerto{It is what makes \enquote{turn the whole picture back by $\alpha$}
  a single rotation applied to both. Turning about two different centres is
  not a rotation of the picture and the cancellation does not happen.}
\item Explain why adding a positional vector to an embedding does not give the
  shift-invariance that rotating it does.
  \answerto{Adding puts the absolute position into the vector, so the dot
  product of two such vectors contains terms in each position separately.
  Rotating puts position into the \emph{angle}, and the identity above removes
  the absolute angles and leaves the difference.}
\item A sequence is shifted bodily along by ten positions. What happens to
  every attention score under a rotary encoding, and why?
  \answerto{Nothing. Every score depends only on a difference of positions,
  and shifting every position by the same amount leaves every difference
  unchanged.}
\item Why does the identity holding at every angle not mean a model works at
  every sequence length?
  \answerto{The mathematics is exact everywhere; whether the model has been
  trained on that part of the circle is a separate, empirical question. It is
  the same distinction as between the cosine being exact and the similarity
  meaning something.}
\item You need to compare directions of vectors whose lengths carry real
  information. Say what cosine similarity would cost you, and name one thing
  you could report alongside it.
  \answerto{It discards the lengths entirely, so two vectors differing only in
  magnitude are indistinguishable. Report the lengths themselves, or use the
  bare dot product, which keeps both.}
\end{furtherproblems}
